\section{The PercentileSweep Algorithm}
\label{sec:algorithm}

\subsection{Formal Definition}

\begin{definition}[\ps{}$(p, T)$]
\label{def:percentile-sweep}
Let $s_0 = \lfloor \text{SHA-256}(\texttt{census\_release\_id} \,\|\,
\texttt{"DIA\_SEED\_V1"}) \rfloor \bmod 2^{31}$ be the content-derived base seed.
For $i = 0, 1, \ldots, T-1$, define the per-index seed:
\[
  s_i = \lfloor \text{SHA-256}(\texttt{census\_id} \,\|\, i) \rfloor \bmod 2^{31}.
\]
For each seed $s_i$, run the \ar{} + METIS pipeline with $s_i$ and record
the normalised edge cut:
\[
  \ECnorm(i) = \frac{\mathrm{EC}(s_i)}{\sqrt{\min(d_1, d_2)}},
\]
where $d_1, d_2$ are the district counts in the two bisection halves.
Sort the $T$ plans by $\ECnorm$ in ascending order (lowest edge cut first):
\[
  \pi_0 \leq \pi_1 \leq \cdots \leq \pi_{T-1}.
\]
Return the plan at rank $r = \floor{p \times T}$, where $p \in [0, 1]$
is the statutory percentile parameter.
The returned plan is $\pi_r$---the plan at the $p$-th quantile of the
$T$-plan bisection distribution.
\end{definition}

The algorithm requires exactly $T$ METIS calls.
Unlike \cs{}, which terminates early when the convergence criterion is
met, \ps{} always runs all $T$ seeds to construct the full distribution.
For $T = 101$, this requires 101 METIS calls per state.
At the observed runtime of approximately 2 seconds per METIS call for
a 50-state congressional map, a complete \ps{} run requires roughly 3.4
minutes per state or about 2.8 hours for all 50 states---comparable to
\cs{}.

\subsection{SHA-256 Seed Derivation}

The seed derivation follows the DIA specification from B.16~\citep{b16paper}.
The base seed is derived from the census release identifier via SHA-256,
ensuring that the seed is unpredictable before census data are published
and deterministically reproducible afterwards.

For \ps{}, seeds are indexed by position $i$ rather than by sequential
advance from $s_0$.
This design choice has two advantages.
First, it allows parallel execution: seeds $s_0, s_1, \ldots, s_{T-1}$
can be computed and executed in any order, with the final ranking applied
after all plans are complete.
Second, it is transparent: any party can independently verify seed $s_i$
for any $i$ without computing all prior seeds.

The per-index seed formula $s_i = \floor{\text{SHA-256}(\texttt{census\_id}
\,\|\, i)} \bmod 2^{31}$ domain-separates the redistricting seed sequence
from other SHA-256 uses of the census identifier~\citep{sha256nist}.
The index $i$ acts as a counter, and the modular reduction ensures
compatibility with 32-bit METIS seed inputs.
The SHA-256 output is truncated to the least-significant 31 bits: bytes 0--3
of the big-endian digest are interpreted as an unsigned 32-bit integer, then
masked to $2^{31} - 1$ (i.e., the high bit is cleared).

\subsection{Relationship to ConvergenceSweep}

\cs{} is the $p = 0.0$ special case of \ps{}.
Both algorithms run a sequence of METIS plans seeded from the census
identifier.
The differences are:
\begin{enumerate}
\item \cs{} uses a sequential stopping rule (halt after $T$ consecutive
      non-improving seeds); \ps{} always runs exactly $T$ seeds.
\item \cs{} returns the global minimum $\ECnorm$; \ps{}$(0.0, T)$ also
      returns the global minimum but does so by sorting and selecting rank 0.
\item \cs{} requires the seeds to be evaluated in order (because the stopping
      criterion depends on the sequence of outcomes); \ps{} seeds are
      independent and parallelisable.
\end{enumerate}

In practice, \cs{} is more efficient for the $p = 0.0$ case because it
terminates early.
For $p > 0$, \ps{} is the appropriate algorithm.

\subsection{The MedianSweep Special Case}

\ms{} = \ps{}$(0.5, T)$ returns the plan at the median of the
$T$-plan bisection distribution.
This is the plan that a uniformly random seed (drawn from the $T$-seed
space) would produce with probability 0.5.

\ms{} has a natural \emph{representativeness} interpretation: it is equally
likely to be drawn by any party running the algorithm with a random seed.
If a legislature wanted to claim that its plan is ``typical'' of what the
algorithm produces rather than optimal, \ms{} provides that posture.

However, the bisection-family median is not the same as the \recom{}-ensemble
median.
The bisection distribution is a distribution over METIS-seeded plans
within the \ar{} family.
The \recom{} ensemble is a Markov-chain distribution over all valid
redistricting plans.
These are different spaces with different coverage properties.

\subsection{The Bisection-Space vs. ReCom-Ensemble Distinction}

This distinction is crucial for interpreting \ps{}.
The $T$ plans generated by \ps{} are all members of the \ar{} bisection
family: they are produced by the same hierarchical recursive bisection
structure, differing only in the leaf-level METIS seed.
They share the same bisection tree topology determined by the prime
factorisation of the district count $k$.

The \recom{} ensemble, by contrast, explores all valid redistricting plans
via random recombination of pairs of adjacent districts.
It is not constrained to the bisection family.
The \recom{} ensemble is much larger and covers a fundamentally different
part of the plan space.

The implication is that \ps{}$(0.5, T)$ is \emph{not} equivalent to
\ts{}(50th, \recom{})---targeting the 50th percentile of the full
\recom{} ensemble.
\ms{} is the median of the bisection family; \ts{}(50th, \recom{}) is
the median of all valid plans.
Section~\ref{sec:legal-posture} discusses the legal implications of this
distinction.

\subsection{Algorithmic Pseudocode}

\begin{algorithm}
\caption{\ps{}$(p, T)$}
\label{alg:percentile-sweep}
\begin{algorithmic}[1]
\Require Statutory parameters $p \in [0,1]$, $T \in \mathbb{Z}_{>0}$
\Require Census identifier \texttt{census\_id}
\State $\text{plans} \leftarrow []$
\For{$i = 0$ to $T - 1$}
  \State $s_i \leftarrow \floor{\text{SHA-256}(\texttt{census\_id} \,\|\, i)} \bmod 2^{31}$
  \State $\pi_i \leftarrow \ar{}(s_i)$  \Comment{Run ApportionRegions with seed $s_i$}
  \State Append $(\pi_i, \ECnorm(s_i))$ to \text{plans}
\EndFor
\State Sort \text{plans} by $\ECnorm$ ascending \Comment{Best (lowest cut) first}
\State $r \leftarrow \floor{p \times T}$
\State \Return $\text{plans}[r].\pi$
\end{algorithmic}
\end{algorithm}

The $T$ iterations in line 2--6 are fully independent and may be
parallelised.
Observed runtime at 50-worker parallelism: approximately $2T/50$
seconds per state for typical census-tract graph sizes.
